\section{Additional discussion}
\label{app:additional-discussion}


\textbf{Related work.} 
While this work and \citet{kobayashi2024wdlowrank} both examine weight decay’s effects on attention matrix rank, this work not only replicates previous results in a larger and more modern setting, but also uncovers novel phenomena. Different from \citet{kobayashi2024wdlowrank}, we find that the rank reduction of attention matrices is associated with better pretraining and fine-tuning performance (up to a weight decay of 0.6 for the OLMo models). In addition, we find that the Query-Key and Value-Projection matrices are differently affected by weight decay, which also has not been reported before.


Prior work has also examined how weight decay affects model behavior. For example, \citet{zhang2025complexity} studies its role in OOD compositional generalization and reasoning-to-memorization transitions under complexity control. \citet{lee2023rethinking} finds that, in self-supervised learning (SSL), weight decay improves the transferability of image representations, but that the optimal value can be difficult to identify because this benefit is not captured by standard SSL evaluation methods. While our work focuses on language model plasticity, the findings from our work—that higher weight decay can improve downstream performance even when it worsens pretraining loss—and prior works are conceptually consistent.

\textbf{Limitations.}
We discuss limitations in more detail below.

\textit{Mechanistic explanations are correlational.} The findings that weight decay promotes more linearly separable representations, reduces attention matrix rank, and reduces overfitting (Section~\ref{sec:explanation}) serve as potential explanations for why weight decay increases plasticity. However, these mechanisms are correlational in nature. Establishing causality here is challenging because it is hard to disentangle mechanisms from other covariates (e.g., changing attention matrix rank while holding model representations constant). Note that, while the mechanisms are correlational, the effect itself—that weight decay improves plasticity (Section 4)—is causal (since the experiments vary only weight decay and keeps other variables constant).


\textit{Results only apply to the scope examined in the paper.}
Our experiments span various model families (Llama-2, OLMo-2), sizes (up to 4B), TPP ratios (20x and 140x), and downstream tasks (six chain-of-thought reasoning tasks, five language understanding and commonsense reasoning tasks, and one safety alignment task). Beyond this scope (e.g., larger models, higher TPP ratios, other fine-tuning tasks, post-training beyond fine-tuning), results may differ. 


\textit{Small sample size for some model groups.}
We pretrain fewer models for some model groups (\llamafourB{} and \olmosevenx{}) due to the large amount of compute required for their pretraining. However, for these models, the results are consistent with other more extensively studied models: 1) better pretraining validation loss does not guarantee better downstream performance and 2) the weight decay that leads to best downstream performance is larger than the standard 0.1 default. 